\chapter{预期研究成果}
\section{方法成果与证据对象}
本课题拟形成一个由标签、输入、模型和评价相互约束的短期尾部预测方法。标签明确公司日及五日结果，输入明确实际公开日期和主体，模型给出收益主干与辅助修正，评价记录年度和随机种子。这些对象共同决定一份风险排序的含义，单独保存网络权重不足以复现研究。

方法成果包括三字段抽取、按公开时点更新的完整关系版本、目标去自身消息及有界残差融合。抽取成果既包含模型也包含独立语义测试；图成果既包含边也包含来源和身份排除记录；融合成果既包含公式和代码，也包含性质、压力和同信息比较。它们分别回答语义表达、关系可信性和修正限制三个问题。

\reporttable{预期成果、使用对象及对应证据}{tab:outputs}{
\begin{tabularx}{\linewidth}{p{3.2cm}YY}\toprule
成果&用途&证据对象\\\midrule
条件尾部标签与四折索引&共同定义模型目标&算法版本、日期边界、分布与独立参照\\
公告规范和抽取模型&从正文获得有限结构&独立原稿、裁决、语义指标与解析失败\\
完整持股版本与主体表&重建预测日关系&原报告、类型/身份、覆盖与年龄\\
有界风险模型&学习受限制的辅助修正&性质推导、梯度检查、同编码器对照\\
正式比较与消融&检验信息和机制增量&全部种子、年度结果、配对区间和失败分析\\
可追溯离线原型&支持有限审核名单研究&输入/模型版本、主干及两路修正、异常记录\\\bottomrule
\end{tabularx}}

\section{风险排序与追溯原型}
图\ref{fig:output}展示拟议输出关系。公司日分数与当日排序名单分开，名单容量按当日可预测公司数计算。结果同时保存主干、公告修正和图修正，让研究人员能够追查高分来自价格状态、事件还是关系。解释这些项时，不把权重和修正方向认作公告的因果作用。

\reportfigure{风险排序、分支修正和版本追溯的输出设计}{fig:output}{\FigOutput}{拟议原型，不含实际公司预测或模拟风险数值。输出为离线研究与审核辅助，不是交易建议。}

原型需能回答一组简单但关键的问题：该行使用哪个模型与索引版本？输入中最晚可用资料是什么？公告是否因解析失败被排除？名单是否存在共享邻居？分数是否经过验证集校准？这些信息不改变AP公式，却决定结果能否被复核。未映射或无辅助资料的公司不应在导出时丢失，而应保留状态并回到相同模型主干。

公开交付以论文、方法代码、配置、数据合同、聚合统计和复现说明为主。授权原始记录、公告正文、逐样本预测、适配器和风险权重按许可留在私有环境。源码包和报告图表不得夹入旧论文正文、个人签署声明或原始数据。

\section{成果判定与论文组织}
论文拟按问题与相关研究、数据时点、方法、实验及总结组织，开题阶段使用学校七栏目描述研究计划。最终主表由冻结预测生成，包含模型、信息权限、测试年、样本量、正类率、AP、差值、区间、每日名单指标及资源消耗。原因/证据抽取与下游排序分表报告，不能用一个较好的抽取分数代替金融预测结果。

如果增量不能稳定成立，研究仍需给出完整实验及失效条件，分析覆盖、抽取误差、年龄分组与市场漂移。成果质量由方法差异、证据完整性和学院评阅决定，不由预先承诺的固定AP增幅决定。是否形成投稿或专利申请，以实际结果及导师判断为准，不将发表等级当作已经获得的成果。
